\section{A Stepanov-type bound for prime counts in Frobenius classes}\label{sec:stepanov}

This section describes the arithmetic heart of the proof. \Cref{thm:cheb-windows} requires, at every
large window scale $x$, a host together with an upper bound with square-root error for the weighted
inert counts $\mathrm{ASK}$ of \cref{sec:hosts}. Such bounds are obtained from a \emph{record}: a list of
real numbers attached to lattices built from chains of theta lifts, subject to eleven systems of
inequalities. Once the record is identified with the arithmetic of the host, the bound follows from
these inequalities by linear arithmetic.

\subsection{Records}

Let $U_0$ be a finite set (later $U_0=A$, the Frobenius classes modulo $A_0$-cosets). For each $u_0\in U_0$
let $W(u_0)$ be a finite non-empty set of \emph{words}; let $\mathrm{Top}$ be a finite set and $m,f\ge1$.

\begin{definition}[Records]\label{def:record}
A \emph{record} $\mathbb D$ is a list of real numbers, indexed as follows:
\begin{itemize}
\item scalars $x,\eps,a,b,b',k,L,L_a,B_0,N_c,\gamma,L_{pL},v,\mathrm{lg},\log\Gamma$;
\item a weight function $a_T$ on $\mathrm{Top}$, and weights $w(u_0,j,\rho)>0$ for $j<m$ and $\rho<f$;
\item for each word $u\in W(u_0)$: a positive weight $d(u_0,u)$, a level $T(u_0,u)>0$, numbers
$\mu_N(u_0,u)\le\mu^{\max}_N(u_0,u)$ and $\mu_M(u_0,u)\le\mu^{\max}_M(u_0,u)$, a number $h(u_0,u)$, and
contributions $E_j(u_0,u)$ for $j<m$;
\item for each $u_0$: numbers $\mu_{CC}(u_0)\ge\mu^{\min}_{CC}(u_0)$, and
$F_o,\mathrm{Ga},\mathrm{Ta},\mathrm{ASK},\mathrm{IO},\mathrm{HH},\delta,\eta,e$;
\item bounds $\bar\delta,\bar\eta,\bar e$.
\end{itemize}
In the records produced from families (\cref{thm:global-chain}), $\mu_N,\mu^{\max}_N$ and
$\mu_M,\mu^{\max}_M$ are the mean and maximal slopes of lattices $N_u$ and $M_u$ attached to the word
$u$, $h$ is the height of a Bessel-type functional, $\mu_{CC}$ is the slope of a lattice assembled from all
words, and $\mathrm{Ga}$ is the divisibility gain at inert places (\cref{rem:stepanov:mechanism}).

Put $\mathrm{ph}(y)=y\log y-y$, $\log\Omega=\gamma+\log(\pi/v)-L_{pL}$ and $E(u_0,u)=\sum_jE_j(u_0,u)$.
Let $s(u_0)=\sum_u d(u_0,u)(\mu_N-h)(u_0,u)/\sum_ud(u_0,u)$ be the weighted mean, and define the
\emph{main term}
\[
\mathrm{Main}(u_0)=\sum_{j<m}\sum_{\tau\in\mathrm{Top}}a_T(\tau)\Big(\log a_T(\tau)-1-\gamma+L_{pL}
-\frac1f\sum_{\rho<f}\log w(u_0,j,\rho)\Big).
\]
\formal{Formalism.RankGapBounds.WeightedMeanData.Data}
\end{definition}

The size conditions $\sfname{SizeData}_2(\mathbb D,C,X'')$, with respect to a list of constants $C$,
are the conjunction of the following eleven systems. They hold for all $u_0$ and all words $u,u'$;
we write $n_T=|\mathrm{Top}|$.
\formal{Formalism.LFunctionChainParameters.LChain.SizeData₂}
\begin{enumerate}[label=(S\arabic*)]
\item\label{S:setting} \emph{Setting}: the obvious positivity conditions, $\bar\delta<1$ and $\bar e\le\frac14$.
\item\label{S:div} \emph{Divisibility}:
$\mu_{CC}(u_0)+h(u_0,u)-F_o(u_0)\le\sum_\tau\big(m(\mathrm{ph}(a_T(\tau))-a_T(\tau)\log\Omega)-a_T(\tau)\log T(u_0,u)\big)
+C_{\rm Div}\big(n_Ta L_a+m(B_0+N_c)x^{e_{\rm Div}\eps}\big)$.
\item\label{S:ecomp} \emph{Comparison of $E$}: $|E(u_0,u)-m\,\mathrm{lg}\sum_\tau a_T(\tau)|\le C_E\,mB_0$.
\item\label{S:hcomp} \emph{Comparison of heights}:
$h(u)+\sum_\tau a_T(\tau)\log T(u)$ varies with $u$ by at most $C_H(n_TaL_a+mB_0)$.
\item\label{S:count} \emph{Count}: $\log T(u_0,u)$ is at least
$\sum_{j}f^{-1}\sum_\rho\log(w(u_0,j,\rho)v/\pi)+m\,\mathrm{lg}-C(L_a+m^2L_a/k)$ and at most $CmL$, and it
varies with $u$ by at most $C\log k$.
\item\label{S:slopes} \emph{Slopes of $M$}: $-C n_TaL_a(1+m/k)\le\mu_M$ and $\mu^{\max}_M\le Cn_TaL_a$.
\item\label{S:quot} \emph{Quotient slopes}:
$\mu_M+\mathrm{Ga}+F_o-E-\mathrm{Lo}\le\mu_N$ and $\mu^{\max}_N\le\mu^{\max}_M+\mathrm{Ga}+F_o-E+C\,mB_0$, where
$\mathrm{Lo}$ is an explicit small term involving $\mathrm{IO}$ and $\log T$.
\item\label{S:loss} \emph{Losses}: $\mathrm{ASK}(u_0)-\mathrm{Ga}(u_0)\le\mathrm{LoP}$ and $\mathrm{IO}(u_0)\le\mathrm{LoP}$,
where $\mathrm{LoP}=C_{\rm Loss}n_Tmb\,x^{e_{\rm Loss}\eps}$.
\item\label{S:slopeineq} \emph{Slope inequality}: if $0\le\delta<1$, $\eta,e\ge0$, $e\le\frac14$ and
$0<s-\mu^{\min}_{CC}\le\mathrm{HH}$, then
\[
s(u_0)-\log|W(u_0)|-\frac{4u_F\log(1+\mathrm{HH}/u_F)+(\eta+4\sqrt e)\,\mathrm{HH}}{1-\delta}\le\mu_{CC}(u_0),
\]
where $u_F=\max_u(\mu^{\max}_N-h)-s$.
\item\label{S:hh} \emph{Height bound}: $s-\mu^{\min}_{CC}\le\mathrm{HH}\le x^{c_{HH}}b$, and $\mathrm{HH}>0$.
\item\label{S:fparams} \emph{Parameters}: $0\le\delta\le\bar\delta$, $0\le\eta\le\bar\eta$ and $0\le e\le\bar e$.
\end{enumerate}

\begin{proposition}[Pointwise bound]\label{prop:stepanov:pointwise}
If $\sfname{SizeData}_2(\mathbb D,C,X'')$ holds, then for every $u_0$
\[
\mathrm{ASK}(u_0)\le\mathrm{Main}(u_0)+\mathrm{Err}_5(u_0),
\]
where $\mathrm{Err}_5$ is an explicit expression in the record and in $C$. Its main contributions are
$\sum_\tau a_T(\tau)\,C(L_a+m^2L_a/k)$, the divisibility error of \ref{S:div}, $\log|W(u_0)|$,
$\mathrm{LoP}$, and $(1-\bar\delta)^{-1}\big(4\log(1+x^{c_{HH}})(\bar U+b)+(\bar\eta+4\sqrt{\bar e})x^{c_{HH}}b\big)$,
where $\bar U$ bounds $u_F$.
\formal{Formalism.LFunctionChainParameters.WeightedSlopeBounds.each_unit_bound_statement}
\end{proposition}

\begin{proof}
By \ref{S:loss}, $\mathrm{ASK}\le\mathrm{Ga}+\mathrm{LoP}$. By the first inequality of \ref{S:quot},
for every word $u$,
\[
\mathrm{Ga}\le\mu_N(u)-\mu_M(u)-F_o+E(u)+\mathrm{Lo}(u).
\]
Take the $d$-weighted mean over $u$. The mean of $\mu_N-h$ is $s(u_0)$. By
\ref{S:slopeineq}--\ref{S:fparams},
\[
s\le\mu_{CC}+\log|W|+(1-\bar\delta)^{-1}\big(4u_F\log(1+\mathrm{HH}/u_F)+(\bar\eta+4\sqrt{\bar e})\mathrm{HH}\big),
\]
where $u_F$ is bounded using the second inequality of \ref{S:quot} and \ref{S:slopes}, and
$\mathrm{HH}\le x^{c_{HH}}b$. Next, \ref{S:div} bounds $\mu_{CC}+h(u)-F_o$; \ref{S:hcomp} allows us to
compare $h(u)+\sum a_T\log T(u)$ across words; \ref{S:ecomp} replaces $E(u)$ by
$m\,\mathrm{lg}\sum a_T$; and \ref{S:slopes} bounds $-\mu_M$. After these substitutions, the
remaining terms
\[
\sum_\tau m\big(\mathrm{ph}(a_T)-a_T\log\Omega\big)-\sum_\tau a_T\log T(u)+m\,\mathrm{lg}\sum_\tau a_T
\]
are bounded by $\mathrm{Main}(u_0)$ plus admissible errors, using the lower bound for $\log T$ in
\ref{S:count}. The terms $\log(v/\pi)$ of \ref{S:count} cancel those of $\log\Omega$.
\end{proof}

\begin{remark}[The Stepanov mechanism]\label{rem:stepanov:mechanism}
The record encodes an inequality between local and global information.
\begin{itemize}
\item The \emph{gain} $\mathrm{Ga}(u_0)$ measures the divisibility of the lattices $N_u$ at the inert
primes: it is a sum over inert places of $\log p$ times the local index gained along the chain of
lifts. By \ref{S:loss} it dominates the inert prime count $\mathrm{ASK}(u_0)$; this is the analogue of
the total order of vanishing of Stepanov's auxiliary function at the rational points.
\item By \ref{S:quot}, this divisibility raises the slope of $N_u$ above that of $M_u$.
\item By the slope inequality \ref{S:slopeineq} and the divisibility bound \ref{S:div}, the slope of $N_u$
is in turn bounded by archimedean data: heights of Bessel-type functionals, the level $T$, and the period
$\log\Omega$. For the hosts used below, $\log\Omega=\gamma+\log(\pi/v)-\Re\,L'/L(1,\chi_{-d})$ is a
Chowla--Selberg quantity \cite{ChowlaSelberg1967}. This is the analogue of the degree bound. The
approximate orthogonality of \cref{thm:theorem-S} enters through the parameters $\delta,\eta,e$ of
\ref{S:slopeineq}, in the manner of Bost's slope inequalities \cite{Bost1996,Bost2001}.
\end{itemize}
The chain length $m$ plays the role of the multiplicity: each inert prime contributes $m$ times to the
gain. The errors in $\mathrm{Err}_5$ are of size $m(B_0+N_c)x^{O(\eps)}+n_Ta L_a(1+m^2/k)+\dots$.
With $m\asymp\sqrt x$, $k\asymp x^{9/2+10\eps}$ and the normalisations of \cref{sec:top}, these errors
are $x^{1/2+O(\eps)}$ relative to the main term (\cref{prop:gap-bound}).
\end{remark}

\subsection{Chain-tied families and the gap bound}

\begin{definition}[Chain-tied families]\label{def:chaintied}
Let $\mathcal F$ be a family with parameter $\eps$, and let $x$ be a window scale.
$\mathcal F$ is \emph{chain-tied} at $x$ with exponents $(e_D,e_L)$ if the following exist:
\begin{itemize}
\item a member $i_0\in I(x)$, whose datum carries a host $H$ on a frame $\varphi$;
\item integers $m,f$, a record $\mathbb D$ with $U_0=A_\varphi$ and $\mathrm{Top}\cong T$, constants
$C,X''$ with $C.e_{\rm Div}\le e_D$ and $C.e_{\rm Loss}\le e_L$, a function $\ell$ compatible with $H$,
an integer $N$, sets $T,T_2\subset A_\varphi$, and functions $a$ and $w$;
\end{itemize}
such that:
\begin{enumerate}[label=(\roman*)]
\item four tie conditions relate these objects to the family; in particular $T$ is the set $T_H$ of
the host;
\item \emph{identification}: $a_T=a$, $w_{\mathbb D}=w$ via the bijections $\mathrm{Top}\cong T$ and
$\{0,\dots,f-1\}\cong\Phi$, $L_{pL}=\Re\ell(\chi_C)$, and
\[
\mathrm{ASK}_{\mathbb D}(u)=\sum_{t\in T}\sum_{j<m}\mathrm{ASK}\big(\mathrm{arith}(H),N,a(t),t,w(u,j,\cdot)\big);
\]
\item $\sfname{SizeData}_2(\mathbb D,C,X'')$ holds, together with error windows, bounds for $B_0$ and
$N_c$, and the size conditions of \cref{thm:cheb-windows}.
\end{enumerate}
\formal{Formalism.MainNotions.ChainTied}
\end{definition}

The identification (ii) ties the record to arithmetic: $\mathrm{ASK}_{\mathbb D}$ is the actual weighted
count of primes inert in $\Q(\sqrt{-d})$, in the Frobenius classes of the host field.

\begin{proposition}[Gap bound]\label{prop:gap-bound}
If $\mathcal F$ is chain-tied at $x$ with exponents $(e_D,e_L)$, then for every $u\in A_\varphi$
\[
\mathrm{gap}(u)\le x^{(e_D+e_L+3)\eps}\,\mathrm{EB},
\]
with $\mathrm{gap}$ and $\mathrm{EB}$ as in \cref{thm:cheb-windows}, provided $\log x\le x^\eps$ and $x$
is large. \formal{Deduction.ch_top}
\end{proposition}

\begin{proof}
By the identification, $\mathrm{gap}(u)=\mathrm{ASK}_{\mathbb D}(u)-\mathrm{Main}(u)$, since
$\mathrm{MAIN}(\ell,a,w)$ is exactly the summand of $\mathrm{Main}$. Apply
\cref{prop:stepanov:pointwise}, and estimate
$\mathrm{Err}_5\le x^{(e_D+e_L+3)\eps}\mathrm{EB}$ using the bounds for $B_0$, $N_c$ and the error
windows.
\end{proof}

\subsection{Links}

Records are produced from families of arithmetic data by \emph{links}. Each link is a property of the
form ``eventually in $x$, for all $i\in I(x)$, an inequality or a structural condition holds for the
datum $D(x,i)$ and its size record''.

\begin{definition}[Links]\label{def:links}
The main links are the following.
\begin{itemize}
\item $\sfname{LinkN}$: for every lift index $j$, a weighted sum over the twisted places $w$ of local
indices $\eta_j(w)$ is at most
\[
d_{Z_0}(j)\Big(C\big(fm\,b\,L^2+(\log x)^2N_c\big)+\sum_\ell\log\ell\sum_{1\le e<t}\mathrm{sh}_\ell(e)\Big)
\]
whenever the relevant Tate--Shafarevich groups are finite. This is the Galois-cohomological link;
the sum of $\mathrm{sh}_\ell$ is bounded by $\sfname{VVle}$.
\item $\sfname{LinkIRR}$: the proportion $\delta_p$ of non-generic rows satisfies
$0\le\delta_p\le\exp(-x^{\eps/(7J)})$.
\item $\sfname{LinkG}$: local lattice lengths $\mathrm{len}_{MN}(u,w)\le d_u(\tau(w)+\iota(w))$ comparing
$M_u$ and $N_u$ at places $w$ not above the level, orbit bounds, and, at inert places of residue
characteristic $p\le X''$, the bounds $0\le\mathrm{gain}(j,w)\le4|T|a/p$ together with a von Mangoldt
class-sum inequality.
\item $\sfname{LinkDIV}$: the divisibility conditions. These are a frame bound, a Bessel identity (the
functionals $B_u$ agree with a fixed pairing up to scalars and isometries), the height bounds
$\mathrm{ht}(B_k(u))\le h(u)$, an upper bound for $\mu_{\rm dual}+h(u)-F_o$ by the main term plus errors,
and the approximate identities behind \ref{S:ecomp} and \ref{S:hcomp}.
\item $\sfname{LinkS}$ (``Theorem~S''): with $\Gamma=\exp(C_\Gamma m^5fL)$, if
$C_0^{-1}\le d_u\Gamma^2/D_*\le C_0$ for all words $u$, then the off-diagonal Hilbert--Schmidt masses
satisfy $\|K_{uu'}\|_{HS}^2\le(1+\eta_0)d_ud_{u'}/D_*$, and the diagonal defects satisfy
$\|K_{uu}-1\|_{HS}^2\le e_0d_u$; moreover $D_{\rm all}\le(1+\eta_1)D_*$. Here
$\eta_0,\eta_1,\sqrt{e_0}\le\exp(-x^{\eps/8})$.
\item $\sfname{LinkMu}$, $\sfname{LinkEta}$ and $\sfname{SmallLevelRankEq}$: slope bookkeeping
and coverage conditions.
\end{itemize}
\formal{Formalism.MainNotions.GlobalTheoremInputs}
\end{definition}

\begin{theorem}[The global theorem]\Fmark\label{thm:global-chain}
Let $\mathcal F$ satisfy the following eleven properties:
\begin{itemize}
\item Siegel-type mass ratios for all members, and the equality of base volumes and masses;
\item explicit two-sided bounds for degrees in terms of gaps;
\item lower and upper slope inequalities;
\item column valuations;
\item free $T$-data, slope lines, Hecke integrality, ramified level zero, and a second-branch
inclusion length.
\end{itemize}
Assume a size condition on $T$. If the seven-share weight condition holds, and $\sfname{LinkG}$,
$\sfname{LinkS}$, $\sfname{LinkIRR}$, $\sfname{LinkDIV}$ and $\sfname{LinkMu}$ hold, then $\mathcal F$
is chain-tied at all sufficiently large window scales.
\formal{Formalism.MainNotions.linkGLOBAL7_min} \formal{Formalism.MainNotions.Lnk.linkCore}
\end{theorem}

\begin{proofoutline}
From the inputs one derives the family-level statements: a count statement (from the mass and volume
properties), slope statements for $M$ and for $N$ (the latter by a lemma on minimal places), lattice
and node statements, frame statements, losses (from the seven-share lemma), characters and windows.
The lemma $\sfname{linkCore}$ then assembles the record. The correspondence is:
\begin{itemize}
\item \ref{S:count} comes from the count statement;
\item \ref{S:slopes} and \ref{S:quot} come from the slope statements, together with $\sfname{LinkG}$;
\item \ref{S:loss} comes from the losses;
\item \ref{S:slopeineq}--\ref{S:fparams} come from the frame statements; the parameter $\delta$ is the
proportion $\delta_p$ of non-generic rows, bounded through $\sfname{LinkIRR}$, and $\eta,e$ come from the
frame statement, which is derived (among other inputs) from $\sfname{LinkS}$;
\item \ref{S:div}--\ref{S:hcomp} come from $\sfname{LinkDIV}$, and \ref{S:hh} from $\sfname{LinkDIV}$ and the
lattice statements.
\end{itemize}
The identification (ii) of \cref{def:chaintied} comes from the characters and the host data. The
slope inequalities rest on exact column-lattice slopes, a short-dual-basis lemma, a base-change formula
for Hermite--Siegel slopes, and a hook-length bound.

%P2PTR thm:global-chain
A complete proof is given in \cref{sec:pf:top-shape}.
\end{proofoutline}

\subsection{The shape theorem}

\begin{theorem}[Shape theorem]\Fmark\label{thm:shape}
If $\sfname{Green}$ holds and $\sfname{HypsAlt}(\sfname{LinkN},\sfname{DefiniteSideCondition})$ is
non-empty, then $\sfname{TopG}$ holds.
\formal{UpperTheory.Top.shape_closed_link_n_of_green}
\formal{GlobalTheory.Fc.forall_goal_of_shape_of_hyps}
\end{theorem}

\begin{proofoutline}
Let $\mathcal A$ be the family class of an element of the hypothesis package. Fix $\lambda>1$ and a
requirement $\mathcal G$ with $\sfname{HostExW}(\lambda,\mathcal G)$ such that the families of
$\mathcal G.\mathrm{cls}(\lambda)$ satisfy $\sfname{VVle}$.
\begin{enumerate}[label=(\alph*)]
\item The field $h_{\rm Ex}$ provides, for all small $\eps$, designed families in
$\mathcal A\cap\mathcal G.\mathrm{cls}(\lambda)$ that are large in the required sense.
\item The links are closed on the relevant class:
  \begin{itemize}
  \item $\sfname{LinkN}$ is a field of the package;
  \item $\sfname{LinkIRR}$ follows from $\sfname{Green}$, through the weights of Gelfand--Graev and
  Jacquet idempotents (\cref{thm:linkIRR});
  \item $\sfname{LinkG}$ follows from a local lattice theorem (\cref{thm:linkG});
  \item $\sfname{LinkDIV}$ follows from bridge-kernel bounds and $\sfname{LinkS}$ from Theorem~S
  (\cref{thm:theorem-S});
  \item the divisibility statement $\sfname{DIVinN}$ follows from $\sfname{LinkN}$, $\sfname{VVle}$ and
  $\sfname{SmallLevelRankEq}$.
  \end{itemize}
\item The mass, volume, slope and integrality properties required by \cref{thm:global-chain} are
fields of the package, or follow from them.
\end{enumerate}
By (a)--(c) and \cref{thm:global-chain}, the designed families are chain-tied at all large window
scales. Hence $\mathcal G.\mathrm{cls}(\lambda)$ is connected. The individual inputs are discussed in
\cref{sec:families}.

%P2PTR thm:shape
A complete proof is given in \cref{sec:pf:top-shape}.
\end{proofoutline}

\begin{remark}
\Cref{thm:shape} is a conditional statement: all of its arithmetic substance lies in the hypothesis
package. This is why the non-emptiness of the package (\cref{thm:hyps-nonempty}) is the largest part
of the formal proof. It accounts for about four million of the eleven million lines in the
dependency closure of the main theorem.
\end{remark}
